% !TEX root = ../matan.tex

\section{Частные производные высших порядков. Первая теорема (Шварца) о равенстве частных производных}

\subsection*{Частные производные высших порядков}

Пусть $f\colon\Omega\to\RR$, $\Omega\subset\RR^d$, и частная производная
$\partial_j f$ определена в окрестности $B(x,\eps)\subset\Omega$. Тогда можно
рассмотреть её частную производную по переменной $x_k$:
\[
\frac{\partial^2 f}{\partial x_k\,\partial x_j}(x)
=\partial_k\partial_j f(x)
:=\frac{\partial}{\partial x_k}\!\left(\frac{\partial f}{\partial x_j}\right)(x).
\]
При $j\ne k$ такая производная называется \textbf{смешанной}. Аналогично определяются
производные более высоких порядков
$\partial_{k_1}\partial_{k_2}\cdots\partial_{k_m} f$.

\begin{Remark}{Порядок дифференцирования важен}{}
	В общем случае результат зависит от порядка дифференцирования:
	\[
	\partial_k\partial_j f \;\ne\; \partial_j\partial_k f.
	\]
	Поэтому равенство смешанных производных требует дополнительных условий
	(непрерывность смешанных производных --- теорема Шварца, либо дифференцируемость
	первых производных --- теорема Юнга).
\end{Remark}

\begin{Example}{Классический контрпример}{}

	Рассмотрим
	\[
	f(x,y)=
	\begin{cases}
		xy\,\dfrac{x^2-y^2}{x^2+y^2}, & (x,y)\ne(0,0),\\[6pt]
		0, & (x,y)=(0,0).
	\end{cases}
	\]
	Прямое вычисление даёт $\partial_x f(0,y)=-y$ при всех $y$ и $\partial_y f(x,0)=x$
	при всех $x$, откуда
	\[
	\partial_y\partial_x f(0,0)=-1,
	\qquad
	\partial_x\partial_y f(0,0)=+1.
	\]
	Смешанные производные в начале координат не совпадают: обе существуют, но не
	являются непрерывными в нуле, и условие теоремы Шварца нарушено.
\end{Example}

\subsection*{Первая теорема о равенстве смешанных производных (теорема Шварца)}

\begin{Theorem}{Шварца}{}
	Пусть $f\colon\Omega\to\RR$, $x$ --- внутренняя точка $\Omega$. Если в некоторой
	окрестности $B(x,\eps)\subset\Omega$ определены смешанные производные
	$\partial_j\partial_k f$ и $\partial_k\partial_j f$, и обе непрерывны в точке $x$,
	то
	\[
	\partial_j\partial_k f(x)=\partial_k\partial_j f(x).
	\]
\end{Theorem}

\begin{proof}
	Без ограничения общности можно считать, что дифференцирование затрагивает лишь
	две переменные, поэтому полагаем $d=2$, $j=1$, $k=2$ (остальные координаты
	фиксированы).

	Введём \emph{приращение по прямоугольнику} с малыми сторонами $\theta_1,\theta_2$:
	\[
	A:=f(x_1+\theta_1,\,x_2+\theta_2)
	- f(x_1+\theta_1,\,x_2)
	- f(x_1,\,x_2+\theta_2)
	+ f(x_1,\,x_2).
	\]
	Величину $A$ мы оценим двумя способами, применяя теорему Лагранжа о среднем в
	разном порядке.

	\textbf{Способ 1} (сначала по $x_2$, затем по $x_1$).
	Зафиксируем $\theta_1$ и рассмотрим вспомогательную функцию
	\[
	\Phi(\tau):=f(x_1+\theta_1,\,\tau)-f(x_1,\,\tau).
	\]
	Тогда $A=\Phi(x_2+\theta_2)-\Phi(x_2)$. Функция $\Phi$ дифференцируема по $\tau$,
	поэтому по теореме Лагранжа существует $t$ между $x_2$ и $x_2+\theta_2$ такое, что
	\[
	A=\theta_2\,\Phi'(t)
	=\theta_2\bigl(\partial_2 f(x_1+\theta_1,t)-\partial_2 f(x_1,t)\bigr).
	\]
	Теперь применим теорему Лагранжа по первой переменной к функции
	$s\mapsto\partial_2 f(s,t)$: существует $s$ между $x_1$ и $x_1+\theta_1$ такое, что
	\[
	\partial_2 f(x_1+\theta_1,t)-\partial_2 f(x_1,t)=\theta_1\,\partial_1\partial_2 f(s,t).
	\]
	Следовательно,
	\[
	A=\theta_1\theta_2\,\partial_1\partial_2 f(s,t),
	\qquad
	\frac{A}{\theta_1\theta_2}=\partial_1\partial_2 f(s,t).
	\]
	При $\theta_1,\theta_2\to 0$ точка $(s,t)\to(x_1,x_2)$, и в силу непрерывности
	$\partial_1\partial_2 f$ в точке $x$
	\[
	\frac{A}{\theta_1\theta_2}\;\xrightarrow{\theta_1,\theta_2\to 0}\;\partial_1\partial_2 f(x).
	\]

	\textbf{Способ 2} (сначала по $x_1$, затем по $x_2$).
	Симметрично, рассматривая
	$\Psi(\sigma):=f(\sigma,x_2+\theta_2)-f(\sigma,x_2)$ и применяя теорему Лагранжа
	сначала по $x_1$, а затем по $x_2$, получаем
	$\tilde s$ между $x_1$ и $x_1+\theta_1$ и $\tilde t$ между $x_2$ и $x_2+\theta_2$
	такие, что
	\[
	A=\theta_1\theta_2\,\partial_2\partial_1 f(\tilde s,\tilde t),
	\qquad
	\frac{A}{\theta_1\theta_2}\;\xrightarrow{\theta_1,\theta_2\to 0}\;\partial_2\partial_1 f(x),
	\]
	где использована непрерывность $\partial_2\partial_1 f$ в точке $x$.

	Левая часть $A/(\theta_1\theta_2)$ в обоих способах одна и та же, поэтому пределы
	совпадают:
	\[
	\partial_1\partial_2 f(x)=\partial_2\partial_1 f(x). \qedhere
	\]
\end{proof}

\begin{Corollary}{Производные высших порядков}{}
	Если все смешанные производные до порядка $m$ непрерывны в окрестности точки $x$,
	то значение $\partial_{k_1}\cdots\partial_{k_m} f(x)$ не зависит от порядка
	индексов $k_1,\dots,k_m$: любая их перестановка даёт тот же результат. Это
	получается последовательным применением теоремы Шварца к соседним индексам
	(транспозиции порождают все перестановки).
\end{Corollary}
